\section{Dataset}

%The following sections describe the data mixture used to create \MixtureVitae{}, along with the pipeline for filtering, processing, and curation.

\vspace{-0.2cm}
We adopt a permissive-first, risk-mitigated strategy, combining sources with clear permissive licenses (e.g. CC-BY, Apache, public domain) with narrowly justified inclusions (government works, EU TDM-eligible data) and targeted synthetic data. Within this framework, the \MixtureVitae{} dataset is constructed from three primary categories:  curated sources for domain-specific expertise, diverse web data for language and general knowledge and instruction-following and reasoning datasets to enhance reasoning and task-completion abilities. 

The major categories of our corpus are visualized in Figure~\ref{fig:domain_comp}. We provide a granular breakdown showing the token count for each component (Figure~\ref{fig:data_comp_deets}), the license distribution (Figure~\ref{fig:license_comp}), and synthetic data usage (Figure~\ref{fig:synthetic_composition}). Specific data sources are detailed in the following subsections.

\begin{figure}[htbp]
    \centering

    % --- Sub-figure 1: Domain Donut Chart ---
    \begin{subfigure}{0.48\textwidth}
        \centering
        \includegraphics[width=\linewidth]{figures/mixture_vitae_category_pie.pdf}
        \caption{Dataset Composition by Top-Level Category and Content Domain}
        \label{fig:domain_comp}
    \end{subfigure}%
    \begin{subfigure}{0.48\textwidth}
        \centering
        \includegraphics[width=\linewidth]{figures/mixture_vitae_license_barchart.pdf}
        \caption{Token Distribution by Governing License}
        \label{fig:license_comp}
    \end{subfigure}

    % --- Main Figure Caption ---
    \caption{Composition of the \MixtureVitae{} dataset (permissive-first, risk-mitigated composition).}
    \label{fig:composition_combined}
\end{figure}


% \begin{figure}[htbp]
%     \centering
%     \includegraphics[width=0.8\linewidth]{figures/synthetic_composition_breakdown.pdf}
%     \caption{Detailed composition of \MixtureVitae{}, broken down by data origin. Each bar represents one of the six primary content domains (matching Figure~\ref{fig:domain_comp}). The stacks visualize the total token count (in billions) for each domain, segmented by source type: \textbf{Non-Synthetic} (human-written permissive text and code), \textbf{Mixed} (sources with partial synthetic data or ambiguous provenance), and \textbf{Synthetic} (data generated by permissive models from permissive seeds).}
% \label{fig:synthetic_composition}
%     \label{fig:syn_data_props}
% \end{figure}

% \begin{figure}[htbp]
%     \centering
%     \includegraphics[width=0.8\linewidth]{figures/license_provenance_tiers.pdf}
%     \caption{Legal provenance and risk-mitigation tiers of the \DatasetSize{}B token \MixtureVitae{} corpus. The dataset is fully segmented into its three constituent legal categories, with all sources falling into a permissive-first or risk-mitigated tier. Token counts (billions) and total corpus percentages are shown for each category.}
% \label{fig:provenance_tiers}
%     \label{fig:syn_data_props}
% \end{figure}

\subsection{Data Sources}
\label{sec:data_sources}
Our dataset selection process is governed by a two-layer  criteria, prioritizing risk mitigation followed by quality and capability objectives:

\textbf{I. Legal \& Licensing:} The primary filter is legal compliance. A dataset is considered only if it operates under a clear permissive license (e.g., CC-BY, Apache 2.0) or is in the public domain. For synthetic data, we further scrutinize the provenance of seed corpora and generator models (Appendix~\ref{app:synth_prov}). The majority of our synthetic sources satisfy full provenance transparency (classified as Tier 1), while a minority of community reasoning datasets with opaque provenance are categorized as Tier 2 to manage residual risk.

\textbf{II. Quality \& Capability:} Among compliant sources, we prioritize datasets with prior evidence of high performance in community mixtures (e.g., \citealp{peS2o}). Furthermore, to address the reasoning deficits typical of strictly permissive web scrapes, we target high-density instruction and reasoning data, a choice driven by the need to boost performance on tasks such as GSM8K~\citep{cobbe2021training} and MMLU\citep{hendrycks2021mmlu}.


% \begin{itemize}
%     \item \textbf{Legal \& Licensing:} The primary filter is legal compliance. A dataset is considered only if it operates under a clear permissive license (e.g., CC-BY, Apache 2.0) or is in the public domain. For synthetic data, we further scrutinize the provenance of seed corpora and generator models (Appendix~\ref{app:synth_prov}). The majority of our synthetic sources satisfy full provenance transparency (classified as Tier 1), while a minority of community reasoning datasets with opaque provenance are categorized as Tier 2 to manage residual risk.
%     \item \textbf{Quality \& Capability:} Among compliant sources, we prioritize datasets with prior evidence of high performance in community mixtures (e.g., \citealp{peS2o}). Furthermore, to address the reasoning deficits typical of strictly permissive web scrapes, we target high-density instruction and reasoning data, a choice driven by the need to boost performance on tasks such as GSM8K~\citep{cobbe2021training} and MMLU\citep{hendrycks2021mmlu}.
% \end{itemize}

The following sections describe each of the three categories of data in \MixtureVitae{}: web, curated sources, and instruction and reasoning datasets.




% \begin{figure}[htbp] 
%     \centering    \includegraphics[width=0.95\linewidth]{figures/mixture_vitae_category_pie.pdf}
%      \caption{Composition of the \MixtureVitae{} dataset, categorized by domain. Percentages represent the proportion of each category within the total 211.2 billion token corpus.}    
%     \label{fig:cat_pie}
% \end{figure}

% \begin{figure}[htbp] 
%     \centering    \includegraphics[width=0.95\linewidth]{figures/mixture_vitae_license_barchart.pdf}
%      \caption{Composition of the \MixtureVitae{} dataset by governing license category, showing total tokens (billions) and percentage of the total corpus. This highlights our risk-mitigated sourcing, demonstrating the dataset is composed entirely of data under clear public domain or permissive licenses. The 'OpenRAIL-M' category corresponds to \textbf{The Stack v1} dataset, and the 'Other Permissive' category consolidates all license types (including MIT and CC-0) that individually constitute less than 2\% of the total corpus.}    
%      % \caption{\textbf{Figure 2:} Composition of the 211.2B token MixtureVitae dataset by governing license category, showing total tokens (billions) and percentage of the total corpus. This highlights our risk-mitigated sourcing, demonstrating the dataset is composed entirely of data under clear public domain or permissive licenses. The 'OpenRAIL-M' category corresponds to The Stack v1 dataset, and the 'Other Permissive' category consolidates all license types (including MIT and CC-0) that individually constitute less than 2\% of the total corpus.}
%     \label{fig:cat_pie}
% \end{figure}


\subsubsection{Web-Scale Corpora}
One subset of our pre-training data is derived from web-scale datasets including Nemotron-CC~\citep{su2025nemotroncc}, MGACorpus~\citep{hao2025reformulationpretrainingdataaugmentation}, and FineFineWeb~\citep{map2024finefineweb}. It also contains synthetic data generated by rephrasing web text from Nemotron-CC and MGACorpus.

% Nemotron-CC and MagaCorpus also contain synthetic data generated by rephrasing web text.

% Nemotron-CC and MagaCorpus also contain synthetic data generated by rephrasing web text.

% The datasets also contain synthetic data generated by rephrasing web text.

\subsubsection{Curated Datasets}
To incorporate domain-specific knowledge and high-quality text, we curate diverse sources: public financial documents from SEC EDGAR~\citep{sec_edgar}, multilingual encyclopedic articles from MegaWika~\citep{barham2023megawikamillionsreportssources} and TxT360~\citep{txt360}, scientific papers from arXiv~\citep{clement2019use} and peS2o~\citep{peS2o}, medical data from Pubmed~\citep{PubMed}, code from The Stack v1~\citep{kocetkov2023the}, patents from the USPTO database~\citep{uspto} and EuroPat~\citep{heafield-etal-2022-europat}, mathematical problems from Deepmind Math~\citep{saxton2019analysingmathematicalreasoningabilities}, and video transcripts from both VALID~\citep{Huu2024VALID} and the YouTube Commons corpus~\citep{Langlais2024}, news and law data from the Open License Corpus~\citep{min2024silo}.

We source 12.6\% of our dataset from \textbf{The Stack v1}, a permissive-first, risk-mitigated code dataset governed by the OpenRAIL-M license. We discuss its permissiveness situation in Appendix~\ref{app:stack}.
%, which, unlike our other sources, provides heuristic license provenance, as it was curated by filtering GitHub repositories for permissive license files. While not guaranteeing every individual file's status, this heuristic represents the largest available permissively filtered code collection essential for coding performance. This "best-effort" curation aligns with our risk-mitigation goals.

% \vspace{-0.2cm}
\subsubsection{Instruction and Reasoning Datasets}
To enhance instruction-following and reasoning, we follow \cite{abdin2024phi4technicalreport} by including considerable synthetic and web-rewrite data. We extensively use fully and partially synthetic data --- all generated from permissive or public-domain seed data using models under permissive licenses.%~\citep{raffel2020exploring,laurencon2022bigscience,lee2022deduplicating,xie2023doremi}.


\textbf{General Instruction Following} We include a strong instruction-following baseline with the Magpie Collection~\citep{xu2024magpiealignmentdatasynthesis}, its derivatives (e.g., Magpie-Phi3-Pro). This is augmented with preference data from UltraFeedback~\citep{cui2024ultrafeedback} and NVIDIA's SFT data blend~\cite{nvidia_sft_datablend_v1}, which contains a curated mixture of permissively licensed subsets from public datasets, including OASST~\citep{köpf2023openassistant}, CodeContests~\citep{Li_2022}, FLAN~\citep{chung2022scalinginstructionfinetunedlanguagemodels}, OpenPlatypus~\citep{platypus2023}, and the training split of GSM8K~\citep{cobbe2021training}. Additionally, we augment the P3~\citep{sanh2021multitask} dataset with a few-shot and multiple-choice format.

\textbf{Reasoning} To improve reasoning, we incorporate general corpora such as Glaive-AI Reasoning Dataset~\citep{glaive_reasoning_2023} and OpenThoughts~\citep{guha2025openthoughts} as well as domain-specific datasets: the legal dataset CaseHOLD~\citep{zheng2021does}, scientific Q\&A from the OpenScience collection~\citep{NVIDIAOpenScience2025}, and agent-focused instructions from OpenManus-RL~\citep{openmanus_rl}.

\textbf{Mathematics and Coding} To strengthen quantitative reasoning, we combine our internally developed synthetic Math Word Problems dataset~(Appendix \ref{app:mwp}) with established datasets like MetaMathQA~\citep{yu2024metamathbootstrapmathematicalquestions} and DM-Math~\citep{saxton2019analysingmathematicalreasoningabilities}, further enriched with large-scale math instruction sets, including OpenMathInstruct-2~\citep{toshniwal2024openmathinstruct118millionmath}, DART-MATH~\citep{tong2024dartmath}, Nemo-Math~\citep{mahabadi2025nemotron}, and Prism-Math~\citep{nvidia2025prismmath}.
 For coding, we combine the Ling Coder collection~\cite{codefuse2025samplemattersleveragingmixtureofexperts} with executable instructions from the StarCoder dataset~\cite{kocetkov2023the} to target a wide range of software engineering tasks.
 
%\paragraph{Web-rewrite} To add high-quality data to the web portion of our dataset with publicly available web-rewrite data that follows the provenance requirements of our organic web data (e.g., cc-by-*, or government data).
% \paragraph{Need for Synthetic Data in Instruction-Following and Reasoning.} Synthetic material is introduced selectively to address domains where permissive human text is scarce—especially multi-step reasoning, instructional tasks, and mathematics. Encyclopedic and curated human sources provide most factual coverage, while synthetic data serves as a focused supplement.


\begin{figure*}[!tb]
\centering
\begin{subfigure}[t]{0.49\textwidth}
    \centering
    \includegraphics[width=\textwidth]{figures/synthetic_composition_breakdown.pdf}
    \caption{\MixtureVitae{} composition by origin (total token counts at the top in billions). Each bar represents one of the six primary content domains (see Fig.~\ref{fig:domain_comp}), segmented by source type: \textbf{Non-Synthetic} (real human-written text and code), \textbf{Mixed} (sources with partial synthetic data), and \textbf{Synthetic} (data generated by permissive models from permissive seeds).}
    \label{fig:synthetic_composition}
\end{subfigure}
\hfill % This adds a flexible space between the two subfigures
\begin{subfigure}[t]{0.49\textwidth}
    \centering
    \includegraphics[width=\textwidth]{figures/license_provenance_tiers.pdf}
    \caption{Legal provenance and risk-mitigation tiers of the \MixtureVitae{} corpus. The dataset is segmented into its three constituent legal categories, with all sources falling into a permissive-first or risk-mitigated tier. Token counts (billions) and total corpus percentages are shown for each category.}
    \label{fig:provenance_tiers}
\end{subfigure}
\caption{Composition and provenance of \MixtureVitae{}: \textbf{(a)} Synthetic-status distribution across the six content domains, \textbf{(b)} licensing tiers and risk posture for the corpus.}
\vspace{-0.5cm}
\label{fig:mv-synth-legal}
\end{figure*}

\vspace{-0.5cm}
\subsubsection{Licensing Tiers and Risk Profiles}
\label{sec:license_tiers}

To make the provenance and legal footing of \MixtureVitae{} transparent, we conceptualize all dataset components into \emph{tiers} based on license type and expected risk profile (see Figure~\ref{fig:provenance_tiers} and Table~\ref{tab:provenance}).\footnote{The high-level groupings presented in this section (e.g., ``Code \& Tech'', ``Reasoning'') and the shard breakdowns in the Appendices are primarily organizational abstractions for visualization and provenance tracking. In practice, the actual training data construction follows a granular \textit{domain-aware mixing strategy} (detailed in Section~\ref{sec:training_curation}), where documents are clustered by base URL or provenance to preserve domain coherence per sample, rather than strictly sampling from rigid high-level partitions.} 

\textbf{Tier 1 --- Explicit Open Licenses \& Public Domain.}  
This tier encompasses text and code under clear permissive licenses (e.g., CC0, CC-BY, Apache 2.0, MIT, BSD, a permissive subset of P3) or in the public domain, such as encyclopedic resources, scientific papers, and portions of curated math corpora. Because licenses are explicit and permissive, the legal risk of reuse is minimal. This tier also includes synthetic data generated from permissively licensed models and seed data. 

\textbf{Tier 2 --- Curated Permissive Corpora with Upstream Opacity.}  

\textbf{(a) Permissive Corpora With Partial or Unverified Provenance.}
    This subset includes resources such as \textsc{The Stack V1} and Wikipedia-derived corpora. The released dataset all carries a permissive license, and curators apply filters (e.g., repository-level license heuristics). However, because provenance is only partially tracked at the file or example level, there remains some residual uncertainty about the licensing status of individual items, hence its separation from Tier~1. This Tier also includes datasets that have no license, but the underlying data is public domain or permissive and requiring the same license as the upstream data, or where the data is solely obtained synthetically from a model that is permissively licensed.
    
    \textbf{(b) Synthetic Data with Non-Permissive or Unverifiable Generators or Seeds.} This tier contains datasets that are themselves permissively licensed (e.g., Apache/MIT/CC-BY), but where either (i) the generator model used to create the synthetic data operates under a more restrictive license (e.g., Llama-3 community license, OpenAI API terms), or (ii) the seed data contains slices whose provenance cannot be fully audited (e.g., partially opaque community mixtures). These datasets constitute only $\approx$4\% of \MixtureVitae{} and are isolated for transparency so that users who require a strictly permissive generator and seed provenance can exclude them (more detail in Table~\ref{tab:provenance}). %We include them in Tier 2(b) rather than Tier 1 because, despite the permissive dataset license, the underlying generation chain involves uncertainty beyond our Tier-1 criteria.
    %where the data is permissively licensed, but the model license is of higher risk, and (3) where the dataset has no license, but the underlying data is public domain or permissive and requiring the same license as the upstream data, or where the data is solely obtained synthetically from a model that is permissively licensed.

\textbf{Tier 3 --- Civic / Governmental Works.}  
This tier includes materials that are either statutory public domain (e.g., U.S. federal works) or under a strong public-purpose rationale for reuse (e.g., government websites, regulatory notices). While not always explicitly licensed, such work---typically created for dissemination---is widely recognized as low-risk for inclusion. Filtering with copyright keyword checks further reduces the possibility of inadvertently including restricted content.

% \paragraph{Note on Dataset Taxonomy vs. Training Mix.}
% The high-level groupings presented in this section (e.g., ``Code \& Tech'', ``Reasoning'') and the shard breakdowns in the Appendices are primarily organizational abstractions for visualization and provenance tracking. In practice, the actual training data construction follows a granular \textit{domain-aware mixing strategy} (detailed in Section~\ref{sec:training_curation}), where documents are clustered by base URL or provenance to preserve domain coherence per sample, rather than strictly sampling from rigid high-level partitions.

\subsection{Data Processing Pipeline}
\label{sec:data_proc_pipe}
To transform the raw data sources into a high-quality and permissively licensed pretraining corpus, we develop a multistage data processing pipeline. Our curation pipeline includes the following stages: ensuring permissive licensing, filtering for CSAM and offensive language, improving overall content quality, and reducing data redundancy. The following sections detail each component.

\label{sec:positive_inclusion}
\textbf{Permissiveness Filtering.} In contrast to standard data pipelines that rely on the retroactive negative filtering of broad web scrapes (e.g., \citealp{fan2025can}), we employ a \textbf{positive inclusion} strategy for web data. Rather than ingesting broad web dumps and filtering post-hoc, we positively select sources based on auditable permissive status. Specifically, we (i) apply an explicit allowlist of governmental and international domains (Appendix~\ref{app:gov_domains}), (ii) curate a set of websites with known permissive licenses (Appendix~\ref{app:permissive_domains}), and (iii) expand this set with risk-mitigated documents by searching for permissive license keywords (e.g., “CC-BY-SA”), excluding documents with restrictive terms (e.g., “all rights reserved”). This upfront design minimizes the risk of including paywalled or opted-out content (e.g., commercial news). We justify the inclusion of governmental works under a strong fair-use rationale, considering their public purpose, content type, and minimal market impact (Appendix~\ref{app:govt_works}).

% \subsubsection{Permissiveness Filtering}
% With our permissive-first, risk-mitigated strategy, we conduct a “pseudo-crawl” of the web corpora, prioritizing permissive sources. Specifically, we (i) apply an explicit allowlist of governmental and international domains (Appendix~\ref{app:gov_domains}), (ii) curate a set of websites with known permissive licenses (Appendix~\ref{app:permissive_domains}), and (iii) expand this set with risk-mitigated documents by searching for permissive license keywords (e.g., “CC-BY-SA”), excluding documents with restrictive terms (e.g., “all rights reserved”). 

\textbf{Quality and Safety Filtering.} Per standard practices~\citep{raffel2020exploring}, we remove documents with base64-encoded text (which can disrupt training) and duplicative headers and footers (e.g., "Home | Search") from FineFineWeb.
We remove obscene, adult and CSAM-related content with keyword-based blocklists adapted from prior work~\citep{NEURIPS2022_ce9e92e3, nakamura2024auroramopensourcecontinual}. For Wikipedia-based documents, we remove articles about films, sporting events, and biographies of living persons in English with applied targeted filtering, to minimize memorization of facts about people, in case of objection to incorrect facts about people being generated by models trained on \MixtureVitae{}.  Besides dataset-level filters, we also evaluate the final model's safety profile via standard red-teaming (Appendix~\ref{app:red_teaming}).

% \textbf{Quality Filtering}

% To reduce redundancy but preserve corpus-scale efficiency, we apply intra-dataset deduplication via prefix-based exact matching, a computationally efficient technique that removes verbatim duplicates by comparing a fixed-length prefix of each document.
\textbf{Deduplication.} Informed by recent findings in large-scale data curation, our deduplication strategy prioritizes diversity over purity. While removing exact repetitions mitigates harmful memorization \citep{lee2022deduplicating}, prior research finds that aggressive, global near-duplicate removal can be detrimental. For example, the creators of the \textbf{FineWeb-Edu} dataset \citep{penedo2024} reported \textit{worsened} model performance by global fuzzy deduplication, postulating that it removed ``too much quality data.'' Therefore, we adopt a local-only approach. We first apply \textbf{intra-dataset deduplication} using prefix-based exact matching to remove verbatim boilerplate text ~\citep{lee2022deduplicating}. We \textbf{intentionally avoid full, cross-dataset fuzzy deduplication} to preserve near-duplicates (e.g., Wikipedia articles with different formatting across sources). We posit that doing so retains \textbf{``stylistic and domain diversity,''} a factor shown to be helpful for model generalization \citep{chen2024diversity}.

%We intentionally avoided enforcing full cross-dataset deduplication. Many sources contain overlapping but non-identical material that can be useful for model training—for example, Wikipedia articles in FineFineWeb, MegaWika, and TxT360 may differ in formatting, metadata, or language coverage. Preserving such near-duplicates helps retain stylistic and domain diversity while still mitigating the harmful effects of exact repetition~\citep{lee2022deduplicating, chen2024diversity_synthetic}.

% [HHN: We can have a note that fineweb decided not to do intra common crawl dump dedups, and that others have shown that duplication of higher quality text is not harmful, but we decided to dedup in order to have a base set that can be used later to upsample.] 

\textbf{Training Example Curation.}
\label{sec:training_curation}
Our process for creating training examples involves several stages: \textbf{1. Heuristic Cleaning.} We remove boilerplate content by eliminating repetitive n-gram prefixes and suffixes, following standard web data cleaning pipelines~\citep{raffel2020exploring}. \textbf{2. Fine-grained Deduplication.} To enhance data quality, we segment documents into sentences and remove duplicate sentences within each document. Documents with high internal repetition (sentence duplication rate $>$ 75\%) are discarded entirely, as this has been shown to improve model performance~\citep{lee2022deduplicating}. \textbf{3. Domain-Aware Mixing.} To construct the final training examples, we employ a domain-aware data mixing strategy~\citep{xie2023doremioptimizingdatamixtures}. Documents are clustered by their base URL (a proxy for domain), and sentences are concatenated first within their original document, then packed with other documents from the same cluster.

% \begin{enumerate}
%     \item \textbf{Heuristic Cleaning:} We remove boilerplate content by eliminating repetitive n-gram prefixes and suffixes, following standard web data cleaning pipelines~\citep{raffel2020exploring}.
%     \item \textbf{Fine-grained Deduplication:} To enhance data quality, we segment documents into sentences and remove duplicate sentences within each document. Documents with high internal repetition (sentence duplication rate $>$ 75\%) are discarded entirely, as this has been shown to improve model performance~\citep{lee2022deduplicating}.
%     % We find, through examination of a sample, that documents with sentence duplicates above 75\% are largely the same documents. 
%     \item \textbf{Domain-Aware Mixing:} To construct the final training examples, we employ a domain-aware data mixing strategy~\citep{xie2023doremioptimizingdatamixtures}. Documents are clustered by their base URL (a proxy for domain), and sentences are concatenated first within their original document, then packed with other documents from the same cluster. 
% \end{enumerate}

\textbf{Additional Filtering for Synthetic Datasets.} To ensure that the synthetic subsets of \MixtureVitae{} adhere to our permissive-first, risk-mitigated approach, we prioritize data originating from seeds that are sourced from permissive sources and generated with models that are themselves permissively licensed. A small portion ($\approx$4\%) of \MixtureVitae{} originates from sources with opaque, mixed or restrictive provenance and is isolated into Tier~2(b), as detailed in Appendix~\ref{app:synth_prov} and Table~\ref{tab:provenance}.

